\section{结论与讨论}

\subsection{主要发现}

我们在官方留一用户协议下，把 JIGSAWS 缝合任务的试次分为新手、中级或专家，遵循本课题规定的实验，并在第一个结果出来之前用书面协议固定了所有模型和比较。简单的试次级运动学统计量明显优于多数类（macro-F1 \FMOne{} 对 \FMZero{}，精确置换 $p = \PMOne{}$）。来自采样帧的视频特征单独使用时没有超过随机水平，无论后期还是早期融合，把它们与运动学结合都没有超过单独的运动学模型；预先声明的主融合模型得分为 \FMThree{}。针对性改进，即参照背景文献思路设计的运动纹理特征，使 macro-F1 基本不变（\FMFour{}）。捷径审计解释了这一模式的大部分。运动学统计量中可迁移的信号主要是试次节奏，它能区分新手与受训外科医生，却不能区分中级与专家；而专家只有两位，模型无法学到两人的共同之处，因此专家类几乎从未被识别出来。整帧外观特征能完美识别外科医生，却不携带可迁移的技能信号；场景动态特征则主要重复了节奏信息。

主要的设计选择取决于问题的规模。只有八位外科医生、每个受训类别只有两位，因此我们使用带强正则化和嵌套选择的线性模型以及冻结的预训练特征，而没有训练深度网络；由于每折只含一个类别时逐折 macro-F1 没有定义，我们报告合并 macro-F1、用户级 bootstrap 区间和精确的用户级置换检验。这些选择让结果如实反映其不确定性，代价是区间较宽，无法分辨小于约 0.1 至 0.2 的 macro-F1 差异。

\subsection{优点与不足}

本研究的主要优点是结论不依赖于看过结果之后做出的选择。折划分是官方的，被留出的外科医生从不参与任何拟合，每个模型只评估一次，负面结果与正面结果一并报告。审计还揭示了分类器实际使用的是哪些信息，这对此类系统的任何临床应用都很重要。第~\ref{sec:differences}~节列出了本方法与背景文献和相关工作的四点不同；对本研究的结论而言，最重要的是以外科医生为单位处理不确定性，以及直接测量身份和节奏捷径。

本方法也有明显的不足。它最擅长的是节奏能揭示的东西，即在标准化台式任务上区分新手和受训外科医生；一旦受训外科医生之间的差别在于质量而非速度，它就很可能失效，专家类的结果正说明了这一点。模型只在七位外科医生上训练，一位不典型的外科医生就可能主导一折；自报技能标签衡量的是练习时长，而不是实际表现。视频特征描述的是整帧画面，没有关注针、线或组织，而观察者恰恰是在这些地方看到运动学无法捕捉的质量差异。最后，等权重后期融合是一种粗糙的组合规则，会让较弱的模态拖累较强的模态。

\subsection{未来方向}

最有效的一步是增加外科医生数量，例如在留出外科医生后用 JIGSAWS 的全部三个任务训练，或加入其他数据集，因为仅凭一位训练外科医生无法学到专家类。区分中级与专家可能需要与节奏关联更弱的特征，例如手势级的时间与错误事件，或按时长归一化的动作质量；本文引入的身份探针和节奏对照可以用来检验任何这类特征。在视频方面，在器械尖端附近区域或缝合最终状态上计算的特征，可能捕捉到组织处理和成品质量，而运动学和我们的整帧描述符都观察不到这些。在训练折内选择融合权重，或采用倾向于较强模态的融合规则，也能避免这里观察到的稀释效应。任何此类扩展都应保持严格的 LOUO 协议和书面预注册，因为外科医生数量少，事后选择很容易过拟合。

\subsection*{AI 工具使用声明}

我们使用了 Claude Code 及 Anthropic 的两个模型 Claude Fable 5.1 和 Claude Opus 5.5，用于规划实验、编写与调试代码、运行实验、生成图表，以及起草和修改本报告与演示文稿。本组审阅了代码、结果和文字，并对所提交工作的内容和质量负责。
